\documentclass[11pt,a4paper]{article}
\usepackage{auditstyle}
\lhead{Fresh Glimm--Jaffe proof audit}
\rhead{16 August 2026}
\newcommand{\eps}{\varepsilon}
\newcommand{\T}{\mathbb T}
\newcommand{\W}{\mathfrak W}
\newcommand{\U}{\mathcal U}
\newcommand{\Tr}{\operatorname{Tr}}

\title{Part II gate scoping in Lamport notation\\[3pt]
\large Minimal logic, momentum-source obstruction, and executed route}
\author{Fresh reconstruction of \texttt{part\_ii/gate\_II11\_scoping.tex}}
\date{16 August 2026}

\begin{document}
\maketitle

\begin{resultbox}
This module records the exact logical gate and the obstruction that dictates the
thermal mean-shift method.  The two transfer lemmas proposed by the scoping document
are not left as promises: their complete proofs appear as Steps
$\langle1\rangle10$ and $\langle1\rangle13$ of the V4 Lamport module.
\end{resultbox}
\tableofcontents

\section{The gate and what is logically sufficient}

For fixed coarse scale $M$, let $\alpha_M^N$ be the D4 lattice refinement and
$\beta_M$ the continuum embedding.  The gate is
\[
 \lim_{N\to\infty}\omega_N(\alpha_M^N(W_M(\xi)))
 =\omega^{\rm ct}_{L,\lambda}(\beta_M(W_M(\xi)))
 \qquad(\xi\in h_{M,L}).
\tag{G}
\]

\LS{1}{1}{Identify the minimum statement that must be proved.}

\LS{2}{1}{It is enough to prove (G) pointwise for every pair $(M,\xi)$.}

\Because The Weyl-reduction module proves that the finite linear span of Weyl
operators is norm dense in each coarse Weyl algebra.  If the state values converge
on each generator, linearity gives convergence on that span.  Given $A$ and an
approximant $A_0$ in the span with $\|A-A_0\|<\delta$, the state norm-one bound gives
\[
 |\omega_N(\alpha_M^N(A))-\rho_M(A)|
 \le2\delta+|\omega_N(\alpha_M^N(A_0))-\rho_M(A_0)|.
\]
Let $N\to\infty$, then $\delta\downarrow0$.  Positivity, normalization, projective
consistency, and continuum identification follow in that module from the same
pointwise limits.

\LS{2}{2}{Uniformity in $\xi$ on bounded sets is therefore not a hypothesis of the
gate.}

\Because Step $\langle2\rangle1$ uses only one fixed approximant, hence finitely
many fixed Weyl symbols.  Quantitative uniform estimates are useful outputs but not
logical prerequisites.

\LS{2}{3}{An $N$-uniform interacting gap is not an input.}

\Because V4 first derives eigenvalue convergence from normal-ordered partition
functions and only then obtains
$E_1(N)-E_0(N)\to\gamma_L>0$.  The eventual gap is an output used in the final
thermal-to-ground interchange.

\LS{2}{4}{The right side of (G) denotes the ground state of the specific $K_L$
constructed as the V3 strong semigroup limit.}

\Because V3 identifies $H_0+V_L$ on the cylinder core but does not prove essential
self-adjointness of that core restriction.  Part II neither needs nor proves
uniqueness among other possible realizations.  The gate is an identification with
the constructed object, in convention $(a,b)=(0,0)$.

\LS{2}{5}{This proves $\langle1\rangle1$.}\QedStep

\section{Why momentum cannot be inserted as a sharp-time Gaussian source}

\LS{1}{2}{Prove the sharp-time momentum-source obstruction.}

Let $0\ne h\in C^\infty(\T_L)$ be real and
$C=(-\partial_t^2-\partial_x^2+m^2)^{-1}$ on
$\R_t\times\T_L$.  The formal momentum insertion would use
$J=\delta_0'\otimes h$.

\LS{2}{1}{The Fourier transform of $J$ is}
\[
 \widehat J(k_0,k)=ik_0\widehat h(k).
\tag{2.1}
\]

\Because Distributional differentiation in time multiplies the Fourier transform by
$ik_0$; the spatial tensor factor supplies $\widehat h(k)$.

\LS{2}{2}{Its covariance norm is}
\[
 \|J\|_{H^{-1}}^2
 =\frac1{2L}\sum_{k\in(\pi/L)\mathbb Z}|\widehat h(k)|^2
 \int_{\R}\frac{dk_0}{2\pi}
 \frac{k_0^2}{k_0^2+m^2+k^2}.
\tag{2.2}
\]

\Because The first-chaos norm of a source is
$\ip J{CJ}$; in Fourier variables, $C$ multiplies by
$(k_0^2+m^2+k^2)^{-1}$.  Insert (2.1) and Parseval's factor $1/(2L)$.

\LS{2}{3}{The right side of (2.2) is infinite.}

\Because Choose $k_*$ with $\widehat h(k_*)\ne0$.  For
$|k_0|\ge\sqrt{m^2+k_*^2}$,
\[
 \frac{k_0^2}{k_0^2+m^2+k_*^2}\ge\frac12.
\]
The corresponding $k_0$ integral dominates the integral of $1/2$ over two
unbounded rays and therefore diverges.  All summands are nonnegative, so one
divergent term makes the sum infinite.

\LS{2}{4}{The same obstruction holds for every spatial lattice.}

\Because Only the spatial momentum set and the value of the positive dispersion
change.  Time remains continuous, and for each nonzero spatial coefficient
$k_0^2/(k_0^2+\gamma_N(k)^2)\to1$ as $|k_0|\to\infty$.

\LS{2}{5}{Hence $\Phi(\delta_0'\otimes h)$ is not a Gaussian first-chaos random
variable and $e^{i\Pi(h)}$ cannot be represented as $e^{i\Phi(J)}$ with an
$H^{-1}$ source.}

\Because A Gaussian source functional exists in $L^2$ precisely when its covariance
norm is finite; (2.2) violates that condition.

\LS{2}{6}{This proves $\langle1\rangle2$.}\QedStep

\LS{1}{3}{Replace the impossible momentum source by the exact jump mean shift.}

\LS{2}{1}{At the free thermal level, compute momentum insertions by the oscillator
trace formula.}

\Because V1 derives the one-mode expression
$\Tr(e^{-tH}D(z))/\Tr e^{-tH}=e^{-|z|^2\coth(t\gamma/2)/2}$ and takes its finite or
trace-dominated infinite product.  No path-space derivative source is introduced.

\LS{2}{2}{At the interacting level, let $e^{i\Pi(p)}$ translate the marked
configuration and shift the Gaussian bridge by the unique harmonic path $H_p$.}

\Because V1 solves
$-\ddot H+\gamma^2H=0$, $H(t)-H(0)=\rho$,
$\dot H(t)=\dot H(0)$ and proves the exact Mehler-chain identity.  The shifted
interaction is the finite Appell sum
\[
 \U(\varphi+H)-\U(\varphi)
 =\lambda\int[4H{:}\varphi^3{:}+6H^2{:}\varphi^2{:}
 +4H^3\varphi+H^4].
\tag{2.3}
\]
Thus the momentum variable enters only through a deterministic, bounded harmonic
interpolant.

\LS{2}{3}{This proves $\langle1\rangle3$.}\QedStep

\section{The executed thermal two-torus route}

\LS{1}{4}{Reduce the gate to eight explicit packages.}

\begin{center}
\begin{tabularx}{0.96\textwidth}{P{1.0cm}Y Y}
\toprule
Item & Required statement & Executed proof \\
\midrule
A1 & Continuum $K_L$, trace class, simple ground, positive gap & V3 Steps
$\langle1\rangle9$--$\langle1\rangle16$ \\
A2 & Lattice and continuum marked-slice mean-shift traces & V1 Step
$\langle1\rangle6$; V3 Step $\langle1\rangle15$ \\
A3 & Exact free thermal forms and fixed-coarse rate & V1 Steps
$\langle1\rangle3$, $\langle1\rangle7$--$\langle1\rangle8$ \\
A4 & Same-noise interaction comparison, Wick dictionary, shifted terms & repaired V2
Steps $\langle1\rangle1$--$\langle1\rangle8$ \\
A5 & Uniform plain and shifted exponential integrability & repaired V2 Step
$\langle1\rangle9$; V3 Steps $\langle1\rangle5$--$\langle1\rangle6$ \\
A6 & Fixed-$t$ thermal convergence & V4 Steps
$\langle1\rangle1$--$\langle1\rangle7$ \\
A7 & Normal-ordered partition convergence, Laplace spectral transfer, derived gap &
V4 Steps $\langle1\rangle8$--$\langle1\rangle11$ \\
A8 & Thermal-to-ground interchange and $C^*$ assembly & V4 Steps
$\langle1\rangle12$--$\langle1\rangle16$ \\
\bottomrule
\end{tabularx}
\end{center}

\LS{2}{1}{A1--A5 imply A6.}

\Because A2 writes both thermal functionals on the common A4 probability space.
A3 compares deterministic free factors; A4 compares all random exponents in $L^2$;
A5 gives uniform integrability; A1 supplies the valid continuum trace identity.
V4 performs the three-term quotient subtraction explicitly.

\LS{2}{2}{A5 and the sourceless part of A6 imply the partition convergence in
A7 after subtracting the lattice zero-point energy.}

\Because The free normal-ordered product converges by the dispersion estimate, and
the interacting factor is $\E e^{-\U_N}$.  The unshifted free product with its
zero-point factor has no nonzero continuum limit; V4 uses only
$\prod_k(1-e^{-t\gamma_N(k)})^{-1}$.

\LS{2}{3}{A7 and the A1 continuum gap imply the uniform spectral-tail input for
A8.}

\Because The exact Laplace-transfer proof gives
$E_1(N)-E_0(N)\to\gamma_L>0$ and
$D_N(t_*)\to D_\infty(t_*)$.  For $t\ge t_*$,
\[
 D_N(t)\le e^{-(t-t_*)[E_1(N)-E_0(N)]}D_N(t_*).
\]

\LS{2}{4}{A6 plus this tail input permits the order of limits to be interchanged,
and the Weyl criterion then proves (G) on the whole inductive-limit algebra.}

\Because V4 Step $\langle1\rangle13$ writes the three-term triangle inequality,
first chooses $t$ uniformly large by the derived gap, and only then sends
$N\to\infty$ using fixed-$t$ convergence.  Its final two steps invoke the proved
Weyl-reduction module.

\LS{2}{5}{This proves $\langle1\rangle4$.}\QedStep

\LS{1}{5}{Explain why the slab-diagonal fallback is not load bearing.}

\LS{2}{1}{A slab approximation at time length $T$ has an error controlled by the
individual lattice gap $\gamma_{N,L}$.}

\Because Spectral expansion gives an excited-state tail with leading factor
$e^{-T\gamma_{N,L}}$.

\LS{2}{2}{Interchanging $T\to\infty$ and $N\to\infty$ on that route would require
an independent eventual uniform lower bound on $\gamma_{N,L}$.}

\Because Without such a bound, no single $T$ makes
$e^{-T\gamma_{N,L}}$ small for all sufficiently large $N$.

\LS{2}{3}{The thermal route is stronger because the trace partition functions
derive precisely that gap through Laplace transfer.}

\Because The compact time circle gives trace identities, and convergence of all
large-time Laplace transforms determines the ordered atoms of the spectrum.  The
slab route has no corresponding normalized trace input.  It is therefore a fallback,
not part of the executed proof chain.

\LS{2}{4}{This proves $\langle1\rangle5$.}\QedStep

\section{Scope conclusion}

The scoping document contributes one substantive negative result, Step
$\langle1\rangle2$, and the dependency architecture in Step
$\langle1\rangle4$.  Its tentative A4 coupling and tentative counterterm were
superseded: the executed coupling is V2's edge-symmetrized same-noise construction,
and the equal-point variance offset cancels the tadpole offset, leaving
$(a,b)=(0,0)$.  The repaired V2 module further supplies the four-edge term and the
fundamental-band/alias partition required for A4 to be true at proof level.

\end{document}
